\section{Evola on Christianity in Europe}

Continuing the translation of the interview in La Nation Europeenne:

\textit{Julius Evola: A Justified Pessimism}\footnote{\url{http://www.geocities.com/Athens/Troy/1856/Muti.htm}}

\textbf{Q}. Do you believe that the influence of Christianity was positive for European civilisation? Don't you think that having adopted a religion of Semitic origin has distorted certain traditional European values?

\textbf{A}. Speaking of Christianity, I often used the expression ``the religion that came to prevail in the West.'' In fact the greatest miracle of Christianity was succeeding in asserting itself among the European peoples, even taking account the decadence into which numerous traditions of these peoples had plunged. Nevertheless, we must not forget the cases in which the Christianisation of the West was only superficial. Besides, if Christianity has, without any doubt, altered certain European values, there are also situations where these values were revived by Christianity, by rectifying and modifying itself. Otherwise, Catholicism would be inconceivable in its various ``Roman'' aspects. In the same way a part of Medieval civilisation would be inconceivable, without phenomena such as the appearance of the great Knightly orders, Thomism, a certain mysticism of a high level (e.g., Meister Eckhart), the spirit of the Crusades, etc.

\paragraph{Comment}

This is an important question and difficult to discuss due to certain subtleties and emotional attachments. On the one hand, the return to a decadent paganism, which can often take the form of a nature mysticism, must be rejected. On the other hand, an ignorant anti-Catholicism will make one miss those admirable phenomena mentioned by Evola.

These phenomena need to be studied from an Evolian point of view, though it is difficult to imagine who would be capable of it. For example, what was the animating spirit of the Knightly orders, particularly the Knights Templar (whose heresy conviction has recently been lifted)? Thomism represented a return to Greek and Roman philosophy. Aquinas pretty much replaced Hebrew law with the Stoic natural law. He also rejected the God of the Hebrews --- the jealous ``old man with a beard''-- in favour of a God more like that of Aristotle and Plato (even the Vedanta according to Guenon). The mystical trend of a Meister Eckhart was based on transcendence and the affirmation of a Self above the human world.

By understanding how these values were ``revived'' in a Medieval Christian form, it will be possible to understand how to revive those values again in the contemporary context.

I would like to add that in the ``Sintesi di dottrina della razza'', Evola writes about ``the great Aryan civilisations of the Orient, to ancient Rome, to the Germano-Roman Middle Ages.'' (p 29)

Out of the ruins, will another arise?


\flrightit{Posted on 2007-12-28 by Cologero}

\begin{center}* * *\end{center}

\begin{footnotesize}
\begin{sffamily}

\texttt{Boreas on 2019-12-28 at 09:38 said:}

A few nights back I saw a dream in which Evola said to me that ``the Truth lies in Jesus''. Maybe Evola-the-harsh-critic-of-Christianity has had some revelations in postmortem states.

\hfill

\end{sffamily}
\end{footnotesize}

